\documentclass[11pt,a4paper]{article}
\usepackage[T1]{fontenc}
\usepackage[utf8]{inputenc}
\usepackage[french]{babel}
\usepackage{lmodern}
\usepackage{amsmath,amssymb}
\usepackage[margin=2.35cm]{geometry}
\usepackage{booktabs,longtable,tabularx}
\usepackage{enumitem}
\usepackage{xcolor}
\usepackage{microtype}
\usepackage{hyperref}

\hypersetup{colorlinks=true,linkcolor=blue!45!black,urlcolor=blue!45!black}
\setlist{nosep}
\newcommand{\code}[1]{\texttt{#1}}
\newcommand{\Kaut}{K^{*}_{\mathrm{aut}}}

\title{Protocole pré-enregistré de la campagne M4.2\\
\large effet de l'exposant de concavité $\gamma$ sur la statistique des
avalanches de faillites}
\author{Programme M4.2 --- dossier \code{m4\_2\_credit\_soc}}
\date{Version 1.1 --- 27 juillet 2026, figée avant tout run scientifique\\
\normalsize (v1 figée le 27 juillet avant le volet 0 ; v1.1 intègre les
corrections de l'audit statistique indépendant du même jour, \emph{avant}
le volet 1 : conditions de confirmation 2, 4, 5, 7 chiffrées, convention
\og coupure hors portée \fg{} implémentée dans le code, bootstrap rendu
obligatoire en confirmation, limites méthodologiques documentées ;
aucun critère n'a été modifié après observation d'un résultat de
campagne --- seuls les 2 runs de coût du volet 0 avaient été lancés)}

\begin{document}
\maketitle

\begin{abstract}
Ce protocole fige, avant l'exécution de tout run de campagne (pilotes
compris), les hypothèses, les grilles, les panels de graines, les fenêtres,
les méthodes d'ajustement et les critères de décision de l'étude M4.2. La
stratégie adaptative autorisée après les pilotes est pré-écrite ici : toute
autre modification devra faire l'objet d'un amendement daté, ajouté avant
les runs concernés, et les analyses effectuées hors protocole seront
étiquetées exploratoires. Le modèle est spécifié dans
\code{report/spec\_m4\_2.pdf} ; le cahier des charges est
\code{prompts/PROMPT\_M4\_2.md}.
\end{abstract}

\tableofcontents

\section{Objet et question}

M4.2 généralise la production de M4B en $F_\gamma(K)=A\,K^\gamma$
($0<\gamma<1$, $A=1$ hors ablation) et dérive les taux d'intérêt des
rendements marginaux $m_\gamma(K)=A\gamma K^{\gamma-1}$ (moyenne
géométrique). La question : \emph{la pente $\tau$ de la distribution des
tailles d'avalanches est-elle modulable par $\gamma$, de façon causalement
intelligible et statistiquement robuste, sans détruire le régime
stationnaire d'accumulation--relaxation ?}

Distinctions imposées par le cahier des charges : un changement de $\tau$
n'est ni un déplacement de la coupure $s_c$, ni une variation du taux de
mortalité ou du rapport de branchement $b$, ni un changement d'échelle du
capital, ni une extinction, une croissance non stationnaire ou une censure.

\section{Hypothèses}

\begin{description}
\item[H1 (principale).] $\tau$ --- exposant de la loi de puissance
tronquée ajustée sur les tailles $s\ge s_{\min}=2$ --- dépend de $\gamma$
de façon reproductible au sens des critères du §\ref{sec:criteres}.
Aucune forme (linéaire, monotone) n'est présupposée.
\item[H0 (nulle).] $\tau$ est institutionnel dans la plage stationnaire
--- fixé par les règles ($k\equiv2$, $\eta=\mathrm{Id}$, cancel+destroy)
comme $b\approx0{,}30$ l'était dans M4B --- et $\gamma$ ne contrôle que
des échelles : $s_c$, taux d'intérêt, démographie, capital.
\item[H2 (secondaire).] $b$ reste épinglé $\approx0{,}30$ pour tout
$\gamma$ (le volume d'appariement, un round par tête, est inchangé).
\item[H3 (secondaire, prédiction dimensionnelle).] La charge d'intérêts
relative croît comme $\gamma$ ($m_\gamma(\Kaut)=\gamma\,\delta/(1-\delta)$) :
durées de vie plus courtes, davantage de racines de liquidité et population
stationnaire $N/\lambda$ plus faible à grand $\gamma$.
\item[H4 (secondaire).] Le taux contractuel médian du réseau final suit
l'ordre de grandeur $\gamma\,\delta/(1-\delta)$.
\end{description}

Ces hypothèses sont falsifiables séparément ; le rapport final rendra un
verdict par hypothèse, en distinguant exploration et confirmation.

\section{Paramètres}

\subsection{Fixes (toute la campagne)}

$\delta=0{,}05$, $\sigma=0{,}25$, $K_0=25$, $A=1$ (sauf bras d'ablation
\ref{sec:ablation}), $\code{pop\_max}=30\,000$, moteur \code{m4\_2-1}
(condensat inscrit dans chaque manifeste), règles institutionnelles de la
réduction minimale (spec §2--3).

\subsection{Variables}

$\gamma$ (levier scientifique), $\lambda\in\{10,30,100\}$ (taille de
système), graine, horizon $T$, et --- uniquement dans le bras d'ablation ---
$A=A_{\mathrm{norm}}(\gamma)$.

\subsection{Panels de graines}

\begin{itemize}
\item exploration (benchmark, pilotes, taille finie, horizon, ablation,
mécanisme) : graines $\{1,2,3\}$ (graine 1 seule pour benchmark, sonde et
mécanisme) ;
\item confirmation : graines $\{11,\dots,15\}$, disjointes de
l'exploration ; extension pré-autorisée à $\{11,\dots,20\}$ si l'écart-type
inter-graines du contraste principal dépasse la moitié de l'effet estimé ;
\item aucun pooling inter-graines : tous les ajustements sont faits par
run, puis agrégés (moyenne $\pm$ écart-type inter-graines ; contrastes
appariés par graine).
\end{itemize}

\section{Plan d'expérience}

Tous les runs de campagne sont exécutés et stockés comme runs Simulation
Lab (décision utilisateur du 27 juillet 2026), étiquetés
\og M4.2 --- \code{<plan>} --- \code{<cellule>} \fg, avec manifeste
\code{manifests/<plan>.json} associant chaque cellule à ses \code{run\_id}
lab. Les 28 figures par run ne sont générées que pour la confirmation et
les cellules centrales ; les runs d'exploration écrivent
\code{snapshot\_every=0} (instantané final seul) et
\code{individual\_every=0}, options sans effet sur les trajectoires
(testé).

\subsection{Volet 0 --- benchmark de coût}

$\gamma\in\{1/3,\,2/3\}$, $\lambda=30$, $T=4000$, graine 1 (2 runs).
Objectif : mesurer le coût réel M4.2 (les 75\,s/run à $T=2000$, $k=2$ de
M4B ne préjugent pas du coût à $\gamma\neq1/2$). Si un run dépasse
30\,minutes, le plan de taille finie est re-budgété par amendement avant
lancement.

\subsection{Volet 1 --- pilotes (exploratoire)}

Grille $\gamma\in\{1/3,\ 0{,}4,\ 1/2,\ 0{,}6,\ 2/3\}$, $\lambda=30$,
$T=4000$, graines $\{1,2,3\}$ (15 runs), plus une sonde $\gamma=0{,}75$,
graine 1 (run exploratoire déclaré, jamais confirmatoire : transitoire
$\Kaut\approx1{,}3\cdot10^{5}$ jugé coûteux a priori). Diagnostics :
statut, stationnarité (§\ref{sec:stationnarite}), échelles numériques,
comptages d'avalanches $s\ge2$, activité du marché, durée des runs.

\subsection{Stratégie adaptative pré-autorisée après pilotes}
\label{sec:adaptatif}

Sur la base des seuls diagnostics ci-dessus (pas des valeurs de $\tau$),
il est pré-autorisé de :
\begin{enumerate}
\item exclure une cellule $\gamma$ non stationnaire, éteinte, censurée ou
numériquement instable (documentée comme résultat négatif) ;
\item doubler l'horizon ($T=8000$) d'une cellule lente (transitoire
mangeant le burn-in) ou pauvre en avalanches ($<100$ tailles $s\ge2$ dans
la fenêtre par défaut) avant de la déclarer indéterminée ;
\item resserrer ou étendre la grille \emph{dans} $[0{,}25,\ 0{,}75]$ ---
par exemple ajouter des points intermédiaires là où un contraste apparaît
--- avec justification écrite datée dans le journal \emph{avant} les runs
ajoutés ;
\item re-budgéter les volets si le benchmark l'exige.
\end{enumerate}
Il est interdit de choisir les cellules selon les valeurs de $\tau$
observées, et toute autre adaptation exige un amendement daté du présent
protocole.

\subsection{Volet 2 --- taille finie (production)}

$\gamma$ retenus après pilotes (a priori les cinq) $\times$
$\lambda\in\{10,\,30,\,100\}$ $\times$ graines $\{1,2,3\}$, avec
$T=8000$ à $\lambda=10$ (statistique d'avalanches) et $T=4000$ à
$\lambda\in\{30,100\}$ ($\approx45$ runs).

\subsection{Volet 3 --- horizon doublé}

$T=8000$, $\lambda=30$, graines $\{1,2,3\}$ sur au moins les deux cellules
$\gamma$ extrêmes retenues et la cellule centrale $\gamma=1/2$ (robustesse
à l'horizon ; sert aussi de contrôle de transitoire à grand $\gamma$).

\subsection{Volet 4 --- confirmation (graines 11--15)}

Contrastes principaux : les cellules $\gamma$ extrêmes retenues contre le
centre $\gamma=1/2$, à $\lambda=30$, $T=4000$, graines $\{11,\dots,15\}$,
exécutées en lots Simulation Lab marqués \code{keep} avec les 28 figures
et les figures agrégées de lot. Si l'effet observé en exploration dépend
de $\lambda$ (signe ou ordre de grandeur différents entre $\lambda=30$ et
$\lambda=100$), les mêmes contrastes sont répétés à $\lambda=100$ en
confirmation. Analyse par contrastes appariés par graine
(cellule $-$ centre), intervalle de Student à 95\,\%.

\subsection{Volet 5 --- ablation de contrôle d'échelle}
\label{sec:ablation}

$K_{\mathrm{ref}}=\Kaut(1/2)=361$ J ;
$A_{\mathrm{norm}}(\gamma)=19^{\,1-2\gamma}$ (spec §7). Cellules :
$\gamma\in\{1/3,\,2/3\}$ (ou les extrêmes retenus) avec
$A=A_{\mathrm{norm}}(\gamma)$, $\lambda=30$, $T=4000$, graines $\{1,2,3\}$.
Si le contraste principal de $\tau$ est confirmé au volet 4, l'ablation du
$\gamma$ le plus contrasté est répétée avec les graines $\{11,\dots,15\}$.
Rapport séparé : effet total de $\gamma$ à $A=1$ ; effet résiduel à
échelle autarcique égalisée. Si l'effet résiduel perd le signe de l'effet
total, la conclusion attribue l'effet à l'échelle, pas à la courbure.

\subsection{Volet 6 --- mécanisme}

$\gamma\in$ grille retenue, $\lambda=30$, $T=4000$, graine 1,
\code{snapshot\_every=50}, \code{individual\_every=5} : instantanés denses
pour les taux contractuels, le levier, les degrés, les durées de vie et
les expositions. Volet descriptif : il n'entre pas dans les critères de
confirmation mais dans l'histoire causale.

\section{Métriques}

\subsection{Primaire}

$\hat\tau$ : exposant de la loi de puissance discrète \emph{tronquée}
$p(s)\propto s^{-\tau}e^{-s/s_c}$, MLE sur les tailles $s\ge s_{\min}=2$
de la fenêtre par défaut (burn-in $T/4$), par run. C'est la métrique du
verdict H1/H0. ($s_{\min}=2$ primaire pour la comparabilité avec la
campagne M4B ; les singletons sont décrits séparément par leur taux.)

\subsection{Secondaires}

Par run et fenêtre : exposant pur $\hat\alpha_{\mathrm{pur}}$ et coupure
$\hat s_c$ ; scan CSN de $s_{\min}\in[1,10]$ (estimateur
\code{fit\_tail\_discrete}, KS minimal) ; log-normale discrète ; LRT
tronquée/pure ($p$ conservatrice $\tfrac12\chi^2_1$) ; Vuong pure/log-normale ;
rapport de branchement $b=1-\sum n_{\mathrm{roots}}/\sum s$ ; fraction
racines/taille ; profondeur (moyenne, max) ; volume (moyenne, max) ;
susceptibilité $\langle s^2\rangle/\langle s\rangle$ ; taux de singletons ;
quantiles $q_{50},q_{90},q_{99}$ et maximum ; scaling de $\hat s_c$,
$s_{\max}$ et $\chi$ avec $\lambda$ ; démographie ($N/\lambda$, taux de
mortalité, âges au décès, parts de causes) ; crédit et réseau (contrats
par tête, dette/capital, taux médian et $q_{90}$ du carnet final,
concentration des principaux) ; fonctionnement de $\eta$ (taux de succès
des rounds, part des fusions, pool/population) ; intégrité (résidu de
bilan, erreurs de carnet).

\subsection{Incertitudes}

Bootstrap par run (rééchantillonnage des tailles avec remise,
$n_{\mathrm{boot}}=200$, percentile 2{,}5--97{,}5\,\%) pour
$(\hat\tau,\hat s_c,\hat s_{\min})$ ; le scan de $s_{\min}$ est refait par
réplicat. Le bootstrap est \textbf{obligatoire} pour l'extraction des
runs confirmatoires (\code{extract\_metrics.py --bootstrap} ; l'analyse
confirmatoire refuse un run sans bloc bootstrap), facultatif en
exploration. Agrégation inter-graines par moyenne $\pm$ écart-type et,
pour les contrastes confirmatoires, intervalle de Student à 95\,\% des
différences appariées par graine.

\section{Méthodes d'ajustement et de comparaison}

Bibliothèque : \code{scripts/lib\_metrics.py} (MLE discrètes normalisées
par zêta de Hurwitz ; support étendu $[s_{\min},\ \max(4\,s_{\max},500)]$
pour la tronquée et la log-normale ; validées sur données synthétiques par
\code{tests/test\_stats.py} : biais $<0{,}05$ sur $\tau$ pur, récupération
de $(\tau,s_c)$, sélection de $s_{\min}$ implanté, préférence LRT/Vuong
pour la famille génératrice). Interdits : conclusion tirée d'une seule
régression log-log ou d'un $R^2$ ; pooling inter-graines.

Conventions d'identifiabilité :
\begin{itemize}
\item ajustement tenté seulement si $\ge100$ tailles $s\ge2$ dans la
fenêtre ; sinon la cellule-fenêtre est \emph{indéterminée} ;
\item coupure déclarée \og hors portée \fg{} si $\hat s_c>10\,s_{\max}$
(implémentée : drapeau \code{s\_c\_out\_of\_range} et champ
\code{tau\_hat} de \code{compare\_laws}) ;
\item l'exposant primaire est celui de la tronquée ; le rapport cite
toujours le couple $(\hat\tau,\hat s_c)$, jamais $\hat\tau$ seul ;
\item le scan de $s_{\min}$ descend à 50 événements de queue
(\code{min\_tail=50}, hérité de l'estimateur M4) : sous le plancher
primaire de 100, sinon le scan serait tronqué aux petits $s_{\min}$ ; ses
résultats sont diagnostiques, jamais primaires.
\end{itemize}

L'exigence du cahier des charges \og estime conjointement $\tau$, $s_c$
et $s_{\min}$ \fg{} est satisfaite de façon \emph{décomposée}, pas par un
MLE joint à trois paramètres : $(\tau,s_c)$ sont estimés conjointement à
$s_{\min}=2$ figé ; $s_{\min}$ est estimé séparément par le scan CSN sur
la loi pure (et re-scanné par réplicat bootstrap). Justification :
l'identifiabilité jointe des trois paramètres dans une loi à coupure est
notoirement instable ; la décomposition est contrôlée par la condition de
confirmation n°4, qui exige la cohérence des deux voies.

\subsection*{Limites méthodologiques assumées (documentées d'avance)}
\begin{itemize}
\item \textbf{Bootstrap iid intra-run} : le rééchantillonnage des tailles
ignore leur dépendance temporelle éventuelle au sein d'un run ; les IC
\emph{par run} peuvent être trop étroits. La variance décisionnelle est
inter-graines (runs indépendants), correctement traitée par les
contrastes appariés.
\item \textbf{LRT légèrement anticonservateur} : la loi pure est
normalisée sur support infini (zêta de Hurwitz), la tronquée et la
log-normale sur support fini $[s_{\min},\max(4s_{\max},500)]$ — les
familles ne sont pas exactement emboîtées. Taux de faux positifs
\og coupure \fg{} mesuré $\approx15\,\%$ (au lieu de 5\,\%) sur données
pures à $n=10^4$ ; la direction favorise la \emph{détection} d'une
coupure, bénin pour l'usage (la condition n°5 ne repose jamais sur le
LRT seul, et le régime attendu — hérité de M4B — est franchement
tronqué).
\item \textbf{Fenêtres emboîtées} : $T/8$, $T/4$, $T/2$ partagent
l'essentiel de leurs données ; leur accord est un contrôle faible. Les
contrôles réellement indépendants sont l'horizon doublé (volet 3) et les
demi-fenêtres disjointes.
\item \textbf{IC de Student à 5 graines} (df$=4$) : normalité
invérifiable, sensibilité à une graine aberrante — d'où la clause
d'extension du panel à 10 graines.
\end{itemize}

\section{Stationnarité, censure et arrêt}
\label{sec:stationnarite}

\begin{itemize}
\item \textbf{Stationnarité} (sur la fenêtre par défaut) : dérive relative
de la population $|\widehat{\mathrm{drift}}|\le4\,\%$ par fenêtre (pente
OLS corrigée de l'autocorrélation intégrée), \emph{et} demi-fenêtres
cohérentes ($|\Delta N/N|\le5\,\%$ et $|\Delta N|\le2$ écarts-types
poolés), \emph{et} mêmes contrôles sur $K_{\mathrm{tot}}$
($|\Delta K/K|\le10\,\%$ entre demi-fenêtres --- l'échelle de capital
converge plus lentement que la démographie à grand $\gamma$). Une cellule
non stationnaire à $T=4000$ est rejouée à $T=8000$ (clause
\ref{sec:adaptatif}) avant exclusion.
\item \textbf{Censure} : \code{status=explosion} (population
$>\code{pop\_max}$) déclenche le diagnostic obligatoire multi-plafonds
($\{30, 60, 120\}\times10^3$) et l'examen des pentes avant toute
conclusion ; un run censuré n'entre jamais dans $\tau(\gamma)$.
Inversement, une trajectoire arrêtée par le plafond n'est jamais déclarée
stationnaire.
\item \textbf{Extinction} : \code{status=extinction} ; documentée avec son
pas. (Dans M4B, l'extinction n'existait qu'à l'amorçage, probabilité
$e^{-\lambda}$.)
\item \textbf{Arrêt de campagne} : budget plafonné à $\sim$200 runs de
campagne et $\sim$6 h de calcul cumulé (8 cœurs, $\le6$ processus) ; toute
extension est un amendement daté.
\end{itemize}

\section{Critères de décision}
\label{sec:criteres}

\subsection{Confirmation de H1 (\og $\tau$ est pilotable par $\gamma$ \fg)}

Toutes les conditions suivantes, sur le(s) contraste(s) principal(aux)
$\gamma$ extrême vs $1/2$ :
\begin{enumerate}
\item signe et ordre de grandeur du contraste $\Delta\hat\tau$ reproduits
sur les graines confirmatoires $\{11,\dots,15\}$ (IC de Student à 95\,\%
des différences appariées excluant zéro) ;
\item persistance en taille : même signe à $\lambda=100$, \emph{et}
($|\Delta\hat\tau(\lambda{=}100)|\ge\tfrac12|\Delta\hat\tau(\lambda{=}30)|$
\emph{ou} IC apparié à $\lambda=100$ excluant zéro) ;
\item stabilité aux fenêtres $T/8$, $T/4$, $T/2$ et à l'horizon doublé
(variations $<$ la moitié de l'effet) ;
\item non-réductibilité au seuil et à la coupure : (a) le contraste des
exposants du scan de $s_{\min}$ ($\hat\alpha_{\mathrm{scan}}$,
\code{fit\_tail\_discrete}) a le même signe que $\Delta\hat\tau$ ;
(b) le contraste subsiste, en signe et à $\ge50\,\%$ de son amplitude,
quand $\tau$ est ré-ajusté par profil de vraisemblance \emph{à coupure
commune} $\bar s_c=\sqrt{\hat s_c^{(1)}\hat s_c^{(2)}}$
(\code{fit\_powerlaw\_cutoff\_fixed\_sc}) dans les deux cellules ;
\item famille défendable dans chaque cellule comparée : LRT tronquée/pure
$p<0{,}05$, \emph{ou} coupure hors portée (convention ci-dessous) avec
Vuong $z>2$ en faveur de la loi de puissance ; et dans aucune cellule la
log-normale n'est strictement préférée (Vuong $z<-2$) ;
\item cellules stationnaires, non censurées, non éteintes ;
\item $|\Delta\hat\tau|>2\sqrt{s_d^2/n_g+\overline{\mathrm{SE}}^2_{\mathrm{boot}}}$,
où $s_d$ est l'écart-type inter-graines des différences appariées, $n_g$
le nombre de graines et $\overline{\mathrm{SE}}_{\mathrm{boot}}$ la
moyenne des écarts-types bootstrap par run de $\hat\tau$ (condition
largement redondante avec la n°1 ; elle garde un rôle si le bootstrap
par run domine) ;
\item l'ablation $A_{\mathrm{norm}}$ attribue à la courbure une part de
l'effet de même signe (sinon la conclusion est requalifiée en effet
d'échelle).
\end{enumerate}

Convention \og coupure hors portée \fg{} (implémentée dans
\code{lib\_metrics.compare\_laws}) : si $\hat s_c>10\,s_{\max}$, la
tronquée dégénère en pure ; le drapeau \code{s\_c\_out\_of\_range} est levé
et $\hat\tau$ (champ \code{tau\_hat}) est lu sur la loi pure. Le rapport
cite toujours le couple $(\hat\tau,\hat s_c)$ ou le drapeau.

\subsection{Infirmation de H1 / soutien de H0}

IC à 95\,\% du contraste apparié $\Delta\hat\tau$ contenant zéro avec une
demi-largeur $\le0{,}15$ (précision suffisante pour exclure un effet de
l'ordre de celui de $\lambda$ dans M4B), \emph{alors que} les grandeurs
d'échelle ($\hat s_c$, taux, démographie, réseau) varient de façon
confirmée avec $\gamma$. Verdict : exposant institutionnel, échelles
pilotables --- et ouverture de la recherche mécanique (§\ref{sec:ouverture}).

\subsection{Indétermination et clause de complétude}

Identifiabilité insuffisante malgré la clause d'horizon doublé, ou
incohérence entre fenêtres/tailles non résolue, ou variance inter-graines
dominante après extension du panel. L'indétermination est un résultat
documenté, pas un échec à masquer.

\textbf{Clause de complétude} : tout contraste qui ne satisfait ni
l'ensemble des conditions de confirmation ni le critère d'infirmation est
rapporté comme \emph{partiellement établi / indéterminé}, en nommant
explicitement la ou les conditions défaillantes. L'arbre de décision est
ainsi partitionnant : confirmé / infirmé / indéterminé-partiel (avec
motifs).

\subsection{Ouverture structurelle (conditionnelle)}
\label{sec:ouverture}

Si H1 est infirmée ou si l'effet est requalifié en échelle, le programme
poursuit --- conformément au cahier des charges --- par l'identification
du verrou mécanique (invariance de la cible $\sqrt{K_\ell K_b}$, plateau
de $b$, densité imposée par $\eta=\mathrm{Id}$, convention de taux,
fragmentation des expositions\dots) \emph{mesuré} sur les volets 2--6,
puis par la formulation et le test d'\emph{une} hypothèse structurelle
minimale (la plus explicative), avec ablation appariée contre la baseline.
Ce volet sera pré-enregistré par amendement daté avant ses runs, une fois
le verrou identifié ; les contraintes sanctuarisées ($k\equiv2$, dérivation
économique des contrats, notations, définition d'avalanche) s'y appliquent.

\section{Infrastructure et traçabilité}

\begin{itemize}
\item Orchestration : \code{scripts/lib\_lab.py} et
\code{scripts/run\_campaign.py} --- un lot Simulation Lab par cellule,
$\le6$ processus simultanés, \code{OMP\_NUM\_THREADS=1} ; reprise par
manifeste (cellule $\to$ \code{run\_id}, statut, condensat moteur).
\item Identification : \code{cell\_id} = SHA-256 tronqué de la
configuration complète $+$ condensat du moteur ; inscrit dans le manifeste
et rappelé dans l'annexe de traçabilité du rapport final (règle : toute
valeur citée est reliable à un \code{run\_id}).
\item Métriques : \code{scripts/extract\_metrics.py} $\to$
\code{results/metrics/<run\_id>.json} ; tables agrégées dans
\code{results/tables/} ; figures régénérables par
\code{scripts/make\_figures.py}.
\item Journal : \code{JOURNAL.md}, entrée datée par volet (chiffres clés,
décisions, amendements).
\item Environnement : Python 3.12.3, numpy 2.4.3, scipy 1.17.1 ;
moteur \code{m4\_2-1} ; machine 8 cœurs / 31 Go.
\end{itemize}

\section*{Engagement de figement}

Ce protocole est figé avant l'exécution du volet 0. Les seuls degrés de
liberté post-pilotes sont ceux du §\ref{sec:adaptatif}. Les critères du
§\ref{sec:criteres} ne seront pas modifiés après observation des
résultats ; toute analyse supplémentaire sera étiquetée exploratoire.

\end{document}
